%% annex/p08/many-valued-and-fuzzy.tex --- Part 8.5 proof material
%% Only annex:p8.matrix.1 compresses an inherited complete source proof.
%% Every other item is a new local specialization, finite decision argument,
%% or derived K3/LP fact and therefore has no source-ratio target.

\begin{annexitem}{annex:p8.matrix.1}{thm:p08-finite-matrix-tableau}{Finite-matrix tableau soundness and completeness}
\begin{proof}
Full expansion terminates.  Give each unexpanded labelled occurrence its
formula complexity.  Expanding it marks that occurrence and adds only labelled
immediate subformulas, so the sum of unexpanded complexities on each branch
strictly decreases.  Every operation has finite arity and at most
\(\lvert\mathsf V\rvert^n\) preimage tuples; hence every root produces a
finite tree.

If \(\Gamma\not\entails_{\mathcal M}B\), choose a countervaluation \(v\).
The root family contains the root labelled by \(v(C)\in D\) for each
\(C\in\Gamma\) and by \(v(B)\notin D\).  It is compatible with \(v\).
Whenever a compatible branch expands
\(c(A_1,\ldots,A_n){:}a\), homomorphic evaluation gives
\(f_c(v(A_1),\ldots,v(A_n))=a\); the corresponding preimage tuple therefore
produces a successor still compatible with \(v\).  Induction through the
expansions preserves such a branch.  It cannot contain two distinct labels for
one atom, and no compatible compound label has empty preimage.  Thus that root
does not close.

Conversely,\linebreak an unclosed fully expanded root has an open saturated branch.
Each labelled atom has a unique label, since otherwise the branch would close.
Assign it that value and assign every unlabelled atom an arbitrary member of
the nonempty set \(\mathsf V\).  Induction on formula complexity now shows
that every branch label \(A{:}a\) satisfies \(v(A)=a\): the atomic case is the
definition, while saturation supplies for a compound a realizing tuple whose
component labels satisfy the induction hypothesis.  Homomorphic evaluation
then gives the compound label.  Applied to the root, this valuation designates
every member of \(\Gamma\) and gives \(B\) its nondesignated root label.
Hence \(\Gamma\not\entails_{\mathcal M}B\).
\end{proof}
\end{annexitem}

\begin{annexitem}{annex:p8.k3.1}{entry:p08-k3}{K3 finite-tableau adequacy}
\begin{proof}
The K3 matrix has a nonempty finite value set and total truth operations.  Its
designated set and the generic root convention are exactly those in the card,
so Theorem~\ref{thm:p08-finite-matrix-tableau} specializes to the stated K3
adequacy equivalence.
\end{proof}
\end{annexitem}

\begin{annexitem}{annex:p8.k3.2}{entry:p08-k3}{Decidability of finite-premise K3 consequence}
\begin{proof}
If the finite input contains \(m\) variables, enumerate its \(3^m\) K3
valuations, recursively compute every premise and the conclusion, and accept
exactly when no valuation designates all premises while leaving the conclusion
nondesignated.  The search terminates and decides the card's consequence
relation.
\end{proof}
\end{annexitem}

\begin{annexitem}{annex:p8.lp.1}{entry:p08-lp}{LP finite-tableau adequacy}
\begin{proof}
The LP matrix has the same finite total operations as K3 and the designated set
fixed by its card.  Substituting those data into
Theorem~\ref{thm:p08-finite-matrix-tableau} gives precisely the stated LP
adequacy equivalence.
\end{proof}
\end{annexitem}

\begin{annexitem}{annex:p8.lp.2}{entry:p08-lp}{Decidability of finite-premise LP consequence}
\begin{proof}
For an input with \(m\) variables, enumerate all \(3^m\) LP valuations and
evaluate formulas recursively.  A countervaluation is exactly one designating
every premise and leaving the conclusion at \(0\).  The finite search decides
whether such a valuation exists.
\end{proof}
\end{annexitem}

\begin{annexitem}{annex:p8.fde.1}{entry:p08-fde}{FDE finite-tableau adequacy}
\begin{proof}
The support-pair FDE matrix has four values, total De Morgan operations, and the
designated set in the card.  Theorem~\ref{thm:p08-finite-matrix-tableau}, with
exactly those data, yields the claimed finite-premise FDE adequacy equivalence.
\end{proof}
\end{annexitem}

\begin{annexitem}{annex:p8.fde.2}{entry:p08-fde}{Finite-premise FDE decision procedure}
\begin{proof}
For \(m\) input variables, enumerate the \(4^m\) support-pair valuations and
evaluate each formula by the displayed operations.  Reject exactly when some
valuation gives positive support to every premise but not to the conclusion.
This finite exhaustive test decides FDE consequence.
\end{proof}
\end{annexitem}

\begin{annexitem}{annex:p8.matrix.2}{thm:p08-k3-lp-facts}{K3 and LP theorem sets and duality}
\begin{proof}
Giving every variable value \(\tfrac12\) in K3 gives every formula that value,
by induction on \(\neg,\land,\lor\).  Since it is nondesignated, K3 has no
tautologies.

For strong Kleene operations, simultaneous induction shows that if a formula
has crisp value \(0\) or \(1\), every replacement of middle-valued variables
by crisp values preserves that result.  Negation exchanges the crisp values;
for conjunction, crisp \(1\) requires both conjuncts crisp \(1\), while crisp
\(0\) supplies a crisp \(0\) conjunct; disjunction is dual.  Thus an LP
valuation making \(A\) nondesignated, hence value \(0\), has a classical
completion falsifying \(A\).  Conversely, every classical valuation is an LP
valuation.  Therefore LP-valid formulas are exactly the classical tautologies.

Finally, a K3 counterexample to
\(\Gamma\entails_{\mathbf K_3}\Delta\) gives every member of \(\Gamma\) value \(1\) and
every member of \(\Delta\) value \(0\) or \(\tfrac12\).  Under negation these
become, respectively, nondesignated LP conclusions and designated LP premises,
so they are exactly a counterexample to
\(\neg\Delta\entails_{\mathsf{LP}}\neg\Gamma\).  Negating the same argument
reverses the correspondence, proving the equivalence.
\end{proof}
\end{annexitem}
